% ================================================================
% V6 — boundary-pair resource theory
% Structural correction of V5:
%   (i) the resource is defined on the OBC/PBC pair generated by one
%       bulk coupling rule, because NHSE is intrinsically a boundary
%       contrast;
%   (ii) the primary scalar is unitary-invariant biorthogonal spectral
%       conditioning, not minimum similarity to a transpose-symmetric
%       matrix;
%   (iii) the thermodynamic resource is the extensive boundary excess
%       of that conditioning.
%
% IMPORTANT DIAGNOSIS CARRIED OVER FROM V4/V5
% A fixed-time, L-uniformly-bounded Julia-channel construction is, by
% its own uniform-boundedness hypothesis, incapable of seeing the one
% thing that defines NHSE: an extensive, boundary-condition-dependent
% effect. Its resource was positive for gauge-trivial Hermitian phase
% chains and insensitive to OBC vs PBC. That is not a gap to shrink
% with more lemmas; it is the wrong object class.
%
% V5 introduced an additional structural problem. Its quantity
%   inf{ log cond(S) : SHS^{-1} is transpose-symmetric }
% is zero on every finite Toeplitz matrix, because every complex
% Toeplitz matrix is unitarily similar to a complex symmetric matrix.
% The OBC Hatano--Nelson matrix is Toeplitz. Hence that definition
% cannot be the desired NHSE resource even before one asks for a lower
% bound. This version removes that obstruction rather than patching it.
%
% DESIGN REQUIREMENTS
%   (R1) extensive in L under OBC for finite-skin-depth NHSE;
%   (R2) O(1) / zero boundary contribution under PBC;
%   (R3) exact unitary invariance and asymptotic invariance under
%        uniformly bounded, boundary-blind local recouplings;
%   (R4) closed HN form;
%   (R5) explicit GBZ correspondence on a regular finite-range class;
%   (R6) direct numerical evaluation without fixed-time approximations.
% ================================================================
\documentclass[reprint,amsmath,amssymb,aps,nofootinbib]{revtex4-2}

\usepackage{graphicx}
\usepackage{bm}
\usepackage{mathtools}
\usepackage{mathrsfs}
\usepackage{amsthm}
\usepackage{bbm}
\usepackage{hyperref}
\usepackage{microtype}
\usepackage{booktabs}
\usepackage{xcolor}

\newtheorem{definition}{Definition}
\newtheorem{theorem}{Theorem}
\newtheorem{proposition}{Proposition}
\newtheorem{lemma}{Lemma}
\newtheorem{corollary}{Corollary}
\newtheorem{remark}{Remark}
\newtheorem{assumption}{Assumption}
\newtheorem{conjecture}{Conjecture}

\newcommand{\Tr}{\operatorname{Tr}}
\newcommand{\id}{\mathbbm{1}}
\newcommand{\GBZ}{\mathrm{GBZ}}
\newcommand{\NHSE}{\mathrm{NHSE}}
\newcommand{\diag}{\operatorname{diag}}
\newcommand{\norm}[1]{\left\lVert #1\right\rVert}
\newcommand{\spec}{\operatorname{spec}}
\newcommand{\cond}{\operatorname{cond}}
\newcommand{\PP}{\mathsf P}
\newcommand{\PhiC}{\Phi}
\newcommand{\Res}{\mathcal R}
\newcommand{\Free}{\mathfrak F}
\newcommand{\gauge}{\mathfrak G}
\newcommand{\kappae}{\kappa_\lambda}

\begin{document}

\preprint{APS/123-QED}

\title{A Boundary-Sensitive Resource Theory of the Non-Hermitian Skin Effect}

\author{Ann Author}
\affiliation{Authors' institution and/or address}
\date{\today}

\begin{abstract}
The non-Hermitian skin effect (NHSE) is intrinsically a statement about the
contrast between open and periodic boundaries. We therefore formulate the
resource on a boundary-paired object rather than on a single finite
Hamiltonian. For a diagonalizable matrix we define the biorthogonal spectral
conditioning
\(\Phi(H)=N^{-1}\sum_\lambda m_\lambda\log\|P_\lambda(H)\|\),
where \(P_\lambda\) is the spectral projector and \(N\) is the Hilbert-space
dimension. The finite-size boundary resource is the positive open-boundary
excess
\(\mathcal R_L=[\Phi(H_L^{\rm OBC})-\Phi(H_L^{\rm PBC})]_+\),
and its density is \(r=\limsup_{L\to\infty}\mathcal R_L/L\).
The construction is invariant under arbitrary unitary relabelings and is
asymptotically unchanged by any boundary-blind similarity family with
uniformly bounded condition number. For the Hatano--Nelson chain we prove
\(\Phi(H_{g,L}^{\rm OBC})=|g|L-\log L+O_g(1)\) while
\(\Phi(H_{g,L}^{\rm PBC})=0\), giving
\(r=|g|\). For a regular class of finite-range translation-invariant
lattice models, the thermodynamic resource density is determined by the
GBZ decay factors through the average \(|\log|\beta||\) of the OBC bulk
branches. The same quantity is directly computable from finite matrices,
without a fixed-time dynamical approximation. The resulting framework is
model-independent at the level of the resource axioms, while the GBZ
formula is stated under explicit regularity assumptions that exclude
scale-free and critical exceptional-point limits.
\end{abstract}

\maketitle

% ================================================================
\section{Introduction}
\label{sec:intro}

A quantum or operator resource theory specifies an object class, a free set,
free transformations, and a quantity that quantifies the asymptotic resource
content \cite{ChitambarGour,LiuYuan,ChiribellaSupermap}. The non-Hermitian skin
effect requires an additional structural constraint: its defining
macroscopic feature is a boundary-condition contrast. The same local bulk
coupling rule produces substantially different OBC and PBC spectra and bulk
eigenvectors in the skin phase \cite{HatanoNelson,YaoWang,OkumaEtAl2020}.
A resource built from a single fixed-boundary object therefore has to explain
how a property that is, operationally, a comparison between two boundary
realizations is encoded.

The present work implements that observation directly. The basic object is
the pair of finite-size truncations \(\{H_L^{\rm OBC},H_L^{\rm PBC}\}_{L\ge L_0}\)
generated by one fixed local bulk rule. The resource is not the absolute
non-normality of the OBC matrix. Instead it is the \,\emph{excess}
non-normality, measured by biorthogonal spectral conditioning, that appears
under OBC relative to PBC. This distinction matters because translationally
invariant PBC matrices can themselves be non-normal, especially in
multiband models. What is special about NHSE is the macroscopic enhancement
that appears when the same bulk rule is opened.

The choice of spectral conditioning is motivated by a robust fact of
non-Hermitian linear algebra: for a simple eigenvalue, the norm of the
spectral projector is the eigenvalue condition number and measures the
non-orthogonality of left and right eigenvectors. In NHSE systems, right and
left skin modes localize at opposite boundaries, producing exponentially
small biorthogonal overlap and exponentially large conditioning under OBC.
This connection has been quantified previously for NHSE models using
condition-number diagnostics and related measures of non-normality
\cite{NakaiEtAl2024,FengEtAl2025}.

This route also removes a specific obstruction in the preceding
similarity-to-reciprocity construction. The quantity
\(
\inf\{\log\cond(S):SHS^{-1}=S^T\}
\)
cannot detect the OBC Hatano--Nelson chain: every finite complex Toeplitz
matrix is unitarily similar to a complex symmetric matrix
\cite{ChienLiuNakazatoTam}. Since the OBC Hatano--Nelson matrix is Toeplitz,
a condition-number-one unitary competitor already exists. Thus a minimum
similarity cost to transpose symmetry is zero on the very example it was
intended to detect. The present construction avoids that basis-dependent
notion entirely.

The main results are as follows. First, the boundary-excess functional is
exactly unitary invariant and changes by at most an \(L\)-independent amount
under uniformly conditioned, boundary-blind similarities. Second, for the
Hatano--Nelson model its OBC value is \(|g|L-\log L+O(1)\), whereas its PBC
value is exactly zero. Third, under a stated regularity hypothesis for
finite-range one-dimensional non-Bloch bands, the thermodynamic density of
this resource is the GBZ average of the spatial decay exponent. Finally,
the resource is directly computable from spectral projectors or left/right
eigenvectors, with numerical stability requirements that are now known to
be essential in NHSE calculations \cite{FengEtAl2025}.

The word ``resource'' is used here in the operator-resource sense. No claim is
made that every non-Hermitian phenomenon is a resource, nor that every form of
NHSE, including scale-free or critical variants, must have a positive
extensive density. Scale-free NHSE provides an explicit reason to separate
finite-skin-depth extensive NHSE from critical boundary localization
\cite{LiSunYangLi}.

% ================================================================
\section{Boundary-paired object class}
\label{sec:objects}

\begin{definition}[Bulk family]
A one-dimensional finite-range translation-invariant bulk rule with \(q\)
internal orbitals per unit cell is specified by
\begin{equation}
H(\beta)=\sum_{r=-R_-}^{R_+}H_r\beta^r,
\qquad H_r\in\mathbb C^{q\times q},
\end{equation}
with fixed hopping ranges \(R_\pm\). For each system size \(L\), the same
bulk rule defines an OBC block-Toeplitz matrix \(H_L^{\rm OBC}\) and a PBC
block-circulant matrix \(H_L^{\rm PBC}\).
\end{definition}

\begin{definition}[Resource object]
The resource object associated with the bulk rule is the asymptotic family
\begin{equation}
\mathscr H=\big\{(H_L^{\rm OBC},H_L^{\rm PBC})\big\}_{L\ge L_0}.
\end{equation}
The two matrices are not treated as independent physical systems: they are
two boundary realizations of the same bulk data \(\{H_r\}\).
\end{definition}

\begin{remark}
This pairing is not an extra physical degree of freedom. It is bookkeeping
for the experimentally meaningful comparison that defines NHSE. A single
OBC matrix does not, by itself, specify which bulk-boundary change should be
called a skin effect rather than generic non-normality, disorder, or an
exceptional-point phenomenon.
\end{remark}

% ================================================================
\section{Biorthogonal spectral conditioning}
\label{sec:conditioning}

For clarity, the main theorems are stated first for diagonalizable matrices.
Extensions to defective points require a pseudospectral or Jordan-chain
regularization and are treated as a separate critical limit rather than
silently included in the same formula.

\begin{definition}[Spectral projector]
Let \(H\in\mathbb C^{N\times N}\) be diagonalizable. For an eigenvalue
\(\lambda\), let \(P_\lambda(H)\) denote the spectral projector onto its
right eigenspace. For simple \(\lambda\), if
\begin{equation}
Hr_\lambda=\lambda r_\lambda,
\qquad
H^\dagger \ell_\lambda=\lambda^*\ell_\lambda,
\end{equation}
then
\begin{equation}
P_\lambda
=\frac{r_\lambda\ell_\lambda^\dagger}
       {\ell_\lambda^\dagger r_\lambda}.
\end{equation}
\end{definition}

\begin{definition}[Finite-size conditioning functional]
For diagonalizable \(H\in\mathbb C^{N\times N}\), define
\begin{equation}
\boxed{
\PhiC(H)
=\frac{1}{N}\sum_{\lambda\in\spec(H)}
 m_\lambda\log\norm{P_\lambda(H)},
}
\label{eq:Phi}
\end{equation}
where \(m_\lambda\) is the algebraic multiplicity. At a defective eigenvalue we
set \(\PhiC(H)=+\infty\) unless a separate regularization is explicitly
specified.
\end{definition}

For a simple eigenvalue, \(\norm{P_\lambda}\ge1\), and with unit-normalized
right and left eigenvectors one obtains
\begin{equation}
\norm{P_\lambda}
=\frac{1}{|\ell_\lambda^\dagger r_\lambda|}
=:\kappa_\lambda.
\label{eq:kappa-simple}
\end{equation}
Thus \(\PhiC\) is the mean logarithmic biorthogonal condition number.
Unlike a condition number of an arbitrarily normalized eigenvector matrix,
Eq.~\eqref{eq:Phi} is invariant under independent rescalings of the
left/right eigenvectors.

\begin{proposition}[Unitary invariance]
\label{prop:unitary}
For every unitary \(U\),
\begin{equation}
\PhiC(UHU^\dagger)=\PhiC(H).
\end{equation}
\end{proposition}

\begin{proof}
Spectral projectors transform as
\(P_\lambda(UHU^\dagger)=UP_\lambda(H)U^\dagger\), and the spectral norm is
unitarily invariant.
\end{proof}

\begin{proposition}[Uniformly conditioned similarity bound]
\label{prop:similarity}
Let \(S\in GL(N,\mathbb C)\) and \(H'=SHS^{-1}\). Then
\begin{equation}
\boxed{
|\PhiC(H')-\PhiC(H)|\le \log\cond(S).
}
\label{eq:Phi-similarity}
\end{equation}
\end{proposition}

\begin{proof}
The spectral projectors obey
\(P_\lambda(H')=SP_\lambda(H)S^{-1}\). Hence
\begin{equation}
\norm{P_\lambda(H')}
\le \cond(S)\norm{P_\lambda(H)}.
\end{equation}
Applying the same inequality to \(H=S^{-1}H'S\) gives the reverse bound.
Averaging the logarithms proves Eq.~\eqref{eq:Phi-similarity}.
\end{proof}

\begin{remark}
Equation~\eqref{eq:Phi-similarity} is the appropriate finite-size statement
for a reversible local recoupling. Crucially, it is an $O(1)$ change when
\(\cond(S)\) is bounded independently of $L$. Consequently such an operation
cannot change an extensive coefficient proportional to $L$.
\end{remark}

% ================================================================
\section{The boundary-excess resource}
\label{sec:resource}

\begin{definition}[Finite-size boundary resource]
For a resource object \(\mathscr H\), define
\begin{equation}
\boxed{
\mathcal R_L(\mathscr H)
:=\Big[\PhiC(H_L^{\rm OBC})-\PhiC(H_L^{\rm PBC})\Big]_+,
}
\label{eq:Rfinite}
\end{equation}
where \([x]_+=\max\{x,0\}\).
\end{definition}

\begin{definition}[Thermodynamic resource density]
The extensive resource density is
\begin{equation}
\boxed{
 r(\mathscr H)
:=\limsup_{L\to\infty}\frac{\mathcal R_L(\mathscr H)}{L}.
}
\label{eq:r-density}
\end{equation}
\end{definition}

\begin{definition}[Free family]
A resource object is free when \(r(\mathscr H)=0\). The finite-size quantity
\(\mathcal R_L\) is retained as the directly observable pre-asymptotic
resource.
\end{definition}

\begin{definition}[Asymptotically free gauge operation]
A family \(\{S_L\}_{L\ge L_0}\) is an admissible boundary-blind gauge if
\begin{equation}
\sup_L\cond(S_L)<\infty,
\end{equation}
and it acts identically on the two boundary realizations:
\begin{equation}
(H_L^{\rm OBC},H_L^{\rm PBC})
\mapsto
(S_LH_L^{\rm OBC}S_L^{-1},S_LH_L^{\rm PBC}S_L^{-1}).
\end{equation}
Global unitary relabelings are included with \(\cond(S_L)=1\). Block-diagonal
on-site recouplings with a uniformly conditioned local block are also included.
\end{definition}

\begin{theorem}[Asymptotic gauge invariance]
\label{thm:gauge}
Under any admissible boundary-blind gauge operation,
\begin{equation}
\boxed{
 r(\mathscr H')=r(\mathscr H).
}
\end{equation}
Moreover, for every finite \(L\),
\begin{equation}
|\mathcal R_L(\mathscr H')-\mathcal R_L(\mathscr H)|
\le 2\log C,
\qquad
C:=\sup_L\cond(S_L).
\label{eq:Rfinite-bound}
\end{equation}
\end{theorem}

\begin{proof}
By Proposition~\ref{prop:similarity}, each of the OBC and PBC values of
\(\PhiC\) changes by at most \(\log C\). Therefore their difference changes by
at most \(2\log C\). The positive-part map is 1-Lipschitz, giving
Eq.~\eqref{eq:Rfinite-bound}. Dividing by $L$ and taking the limsup proves the
invariance of $r$.
\end{proof}

\begin{remark}
The resource is deliberately a thermodynamic rate. A bounded local operation
may change finite-size numbers by an additive constant without changing the
extensive NHSE content. This is the same structural logic by which many
thermodynamic resource quantities are defined only after regularization.
\end{remark}

\begin{proposition}[Free-family stability]
If \(r(\mathscr H)=0\), then every admissible boundary-blind gauge transform of
\(\mathscr H\) is also free.
\end{proposition}

\begin{proof}
Immediate from Theorem~\ref{thm:gauge}.
\end{proof}

% ================================================================
\section{Hatano--Nelson theorem: exact OBC/PBC separation}
\label{sec:HN}

Consider the open chain
\begin{equation}
H_{g,L}^{\rm OBC}
=m\id+Je^gT_++Je^{-g}T_-,
\qquad T_-=T_+^T,
\label{eq:HN}
\end{equation}
with \(g\in\mathbb R\). For definiteness take $g>0$ below; the result for
negative $g$ follows by spatial reflection.

Let
\begin{equation}
 k_j=\frac{\pi j}{L+1},
\qquad
 s_j(n)=\sin(k_j n),
\qquad j=1,\ldots,L.
\end{equation}
The right and left eigenvectors can be written as
\begin{equation}
 r_j(n)=\frac{e^{g(n-1)}s_j(n)}{R_j},
 \qquad
 \ell_j(n)=\frac{e^{-g(n-1)}s_j(n)}{L_j},
\end{equation}
where
\begin{equation}
R_j^2=\sum_{n=1}^L e^{2g(n-1)}s_j(n)^2,
\qquad
L_j^2=\sum_{n=1}^L e^{-2g(n-1)}s_j(n)^2.
\end{equation}
Since
\(\sum_{n=1}^Ls_j(n)^2=(L+1)/2\), the simple-eigenvalue condition number is
\begin{equation}
\kappa_j
=\frac{2R_jL_j}{L+1}.
\label{eq:HN-kappa}
\end{equation}

\begin{theorem}[Hatano--Nelson OBC scaling]
\label{thm:HN-OBC}
For every fixed $g\ne0$,
\begin{equation}
\boxed{
\PhiC(H_{g,L}^{\rm OBC})
=|g|(L-1)-\log(L+1)+O_g(1),
}
\label{eq:HN-asymptotic}
\end{equation}
and therefore
\begin{equation}
\boxed{
\lim_{L\to\infty}\frac{\PhiC(H_{g,L}^{\rm OBC})}{L}=|g|.
}
\end{equation}
\end{theorem}

\begin{proof}
We again assume $g>0$. The pointwise lower bounds
\begin{equation}
R_j\ge e^{g(L-1)}|\sin k_j|,
\qquad
L_j\ge |\sin k_j|
\end{equation}
combined with Eq.~\eqref{eq:HN-kappa} give
\begin{equation}
\log\kappa_j
\ge g(L-1)+\log\frac{2}{L+1}+2\log|\sin k_j|.
\end{equation}
Using the classical product identity
\begin{equation}
\prod_{j=1}^L\sin\frac{\pi j}{L+1}
=\frac{L+1}{2^L},
\end{equation}
we obtain
\begin{equation}
\PhiC(H_{g,L}^{\rm OBC})
\ge g(L-1)-\log(L+1)-\log2
+\frac{2\log(L+1)}{L}.
\label{eq:HN-lower}
\end{equation}
For the upper bound, use $s_j(n)^2\le1$ to obtain
\begin{equation}
R_j^2\le\sum_{n=1}^Le^{2g(n-1)}
\le\frac{e^{2g(L-1)}}{1-e^{-2g}},
\qquad
L_j^2\le\sum_{n=1}^Le^{-2g(n-1)}
\le\frac{1}{1-e^{-2g}}.
\end{equation}
Thus
\begin{equation}
\PhiC(H_{g,L}^{\rm OBC})
\le g(L-1)-\log(L+1)
+\log\frac{2}{1-e^{-2g}}.
\label{eq:HN-upper}
\end{equation}
Equations~\eqref{eq:HN-lower} and \eqref{eq:HN-upper} prove
Eq.~\eqref{eq:HN-asymptotic}. Replacing $g$ by $|g|$ follows by reflection.
\end{proof}

\begin{theorem}[Hatano--Nelson PBC]
\label{thm:HN-PBC}
For the same bulk hopping rule with periodic closure,
\begin{equation}
\boxed{
\PhiC(H_{g,L}^{\rm PBC})=0
\quad\text{for every }L,g.
}
\end{equation}
Consequently,
\begin{equation}
\boxed{
\mathcal R_L(\mathscr H_g)
=|g|L-\log L+O_g(1),
\qquad
r(\mathscr H_g)=|g|.
}
\label{eq:HN-resource}
\end{equation}
\end{theorem}

\begin{proof}
The PBC matrix is circulant and is diagonalized by the discrete Fourier matrix
$F$. The Fourier eigenvectors are orthonormal, so every spectral projector has
norm one. Hence \(\PhiC=0\). The resource formula then follows from
Theorem~\ref{thm:HN-OBC}.
\end{proof}

\begin{remark}[Physical interpretation]
The value $|g|$ is the inverse skin depth in lattice units for the
exponential gauge family. The theorem does not identify the NHSE with a
mere choice of basis: it identifies the extensive resource with the
boundary-induced collapse of left/right biorthogonal overlap. In particular,
unitary relabelings cannot remove it.
\end{remark}

% ================================================================
\section{General finite-range chains and the GBZ}
\label{sec:GBZ}

We now connect the resource to non-Bloch band theory. The purpose of this
section is not to hide the needed hypotheses behind the word ``universal.''
The statement below is a theorem for a regular class and identifies exactly
where critical models require separate treatment.

\begin{assumption}[Regular GBZ branch]
\label{ass:GBZ}
For a finite-range block symbol $H(\beta)$ with $q$ orbitals per unit cell,
assume:
\begin{enumerate}
\item the bulk OBC spectrum admits $q$ regular GBZ branches
$\beta_\alpha(\theta)$, $\theta\in[0,2\pi)$, with no positive-measure
exceptional or root-coalescence set;
\item on branch $\alpha$ the selected OBC roots have a common modulus
$\rho_\alpha(\theta)=|\beta_\alpha(\theta)|$ that controls the bulk spatial
envelope, while all competing roots are separated by a fixed exponential
gap except at a measure-zero set;
\item the normalized right and left finite-size bulk eigenvectors have
opposite envelope exponents
$\pm\log\rho_\alpha(\theta)$ up to subexponential prefactors;
\item the PBC Bloch blocks have fixed internal dimension $q$; any Bloch
exceptional points are isolated and of finite algebraic order, and the
discrete PBC momenta used at finite $L$ do not coincide exactly with a
defective Bloch point (equivalently, a fixed generic boundary twist may be
used).
\end{enumerate}
\end{assumption}

Assumption~\ref{ass:GBZ} is the standard finite-skin-depth, noncritical regime.
It excludes scale-free NHSE, for which the relevant localization exponent may
vanish with $L$, and it excludes finely tuned extensive exceptional-point
singularities. This exclusion is necessary if the resource is required to
scale linearly with $L$.

\begin{lemma}[GBZ envelope implies conditioning exponent]
\label{lem:GBZ-envelope}
Under Assumption~\ref{ass:GBZ}, a bulk OBC eigenmode on branch $\alpha$ at
parameter $\theta$ satisfies
\begin{equation}
\log\norm{P_{\alpha,\theta}^{(L)}}
=L\,\big|\log\rho_\alpha(\theta)\big|+o(L).
\label{eq:projector-GBZ}
\end{equation}
\end{lemma}

\begin{proof}
The right eigenvector has a spatial envelope proportional to
$\rho_\alpha^n$ up to subexponential factors, while the left eigenvector has
the opposite envelope $\rho_\alpha^{-n}$. If $\rho_\alpha>1$, the right norm
is dominated by the right boundary and the left norm by the left boundary;
for $\rho_\alpha<1$ the roles are reversed. Their unnormalized overlap is a
sum of subexponential terms and therefore contributes only $o(L)$ to the
logarithm. With unit-normalized left/right vectors,
$\norm{P}=1/|\ell^\dagger r|$, giving Eq.~\eqref{eq:projector-GBZ}.
The exceptional set excluded in Assumption~\ref{ass:GBZ} cannot alter the
thermodynamic average.
\end{proof}

\begin{theorem}[GBZ-conditioning correspondence]
\label{thm:GBZ}
Under Assumption~\ref{ass:GBZ},
\begin{equation}
\boxed{
\lim_{L\to\infty}\frac{\PhiC(H_L^{\rm OBC})}{L}
=\frac{1}{q}\sum_{\alpha=1}^{q}
\int_0^{2\pi}\frac{d\theta}{2\pi}
\big|\log\rho_\alpha(\theta)\big|,
}
\label{eq:GBZ-main}
\end{equation}
while
\begin{equation}
\boxed{
\PhiC(H_L^{\rm PBC})=o(L).
}
\label{eq:PBC-subextensive}
\end{equation}
Consequently,
\begin{equation}
\boxed{
 r(\mathscr H)
=\frac{1}{q}\sum_{\alpha=1}^{q}
\int_0^{2\pi}\frac{d\theta}{2\pi}
\big|\log\rho_\alpha(\theta)\big|.
}
\label{eq:GBZ-resource}
\end{equation}
\end{theorem}

\begin{proof}
By Lemma~\ref{lem:GBZ-envelope}, each OBC bulk mode contributes
$L|\log\rho_\alpha(\theta)|+o(L)$ to the logarithm of its spectral-projector
norm. There are $L+o(L)$ modes on each regular branch, and the finite-size
OBC quantization becomes an integral over $\theta$. Dividing the sum in
Eq.~\eqref{eq:Phi} by the total dimension $qL$ gives
Eq.~\eqref{eq:GBZ-main}.

For PBC, Fourier transformation is unitary and reduces the matrix to $L$
independent $q\times q$ Bloch blocks $H(e^{ik})$. Away from isolated exceptional
momenta, each block has finite-dimensional spectral projectors with norms
independent of $L$. Isolated finite-order exceptional points can generate at
most polynomial finite-size divergences after discretization, whose logarithm
is $o(L)$. Thus Eq.~\eqref{eq:PBC-subextensive} follows, and inserting it into
Eq.~\eqref{eq:Rfinite} yields Eq.~\eqref{eq:GBZ-resource}.
\end{proof}

\begin{corollary}[Finite-skin-depth criterion]
\label{cor:GBZ-criterion}
Within the regular class of Theorem~\ref{thm:GBZ}, the boundary resource density
vanishes if and only if
\begin{equation}
\rho_\alpha(\theta)=1
\quad\text{for almost every }(\alpha,\theta).
\end{equation}
Thus the resource detects the GBZ departure from the unit circle in precisely
the finite-skin-depth regime.
\end{corollary}

\begin{remark}[Scope of the word ``universal'']
The construction of \(\Res\) is model-independent over the boundary-paired
lattice class. The closed GBZ formula is universal only within the regular
finite-skin-depth class of Assumption~\ref{ass:GBZ}. Scale-free NHSE and
critical exceptional-point limits require a different scaling functional;
forcing them into a linear-in-$L$ resource would contradict their defining
finite-size scaling.
\end{remark}

% ================================================================
\section{Resource-theoretic interpretation and physical free operations}
\label{sec:free}

The free transformations used here are deliberately narrower than arbitrary
nonunitary similarities. Their physical role is that of boundary-blind changes
of representation or local reciprocal recoupling, not a license to modify the
bulk nonreciprocity itself.

\subsection{Global unitary relabeling}
A change of Hilbert-space coordinates
\begin{equation}
H_L\mapsto U_LH_LU_L^\dagger
\end{equation}
is physically free. Proposition~\ref{prop:unitary} gives exact invariance of
both OBC and PBC conditioning and hence of \(\mathcal R_L\) and $r$.
This removes the basis dependence that invalidated the previous
similarity-to-transpose-symmetry resource.

\subsection{Local reciprocal recoupling}
Consider a local on-site basis change with a block-diagonal matrix
\begin{equation}
S_L=s_1\oplus s_2\oplus\cdots\oplus s_L,
\qquad
\sup_n\cond(s_n)\le C_0<\infty.
\end{equation}
Then
\begin{equation}
\cond(S_L)=\max_n\cond(s_n)\le C_0,
\end{equation}
so Theorem~\ref{thm:gauge} applies. The leading extensive coefficient of the
boundary resource is unchanged. This is the precise operational content of
calling such a recoupling free: it cannot generate an $O(L)$ boundary excess
from an $O(1)$ one.

\subsection{What is not free}
A site-dependent nonunitary gauge whose total condition number grows as
$e^{cL}$ is not an asymptotically free transformation. Such a transformation
carries an extensive normalization cost, exactly as an extensive resource
should. This distinction is essential in the Hatano--Nelson problem: the
similarity that removes asymmetric hopping under OBC has condition number
$e^{|g|(L-1)}$, and its exponential growth is the same scale that appears in
the biorthogonal conditioning resource.

\begin{remark}[On conversion theory]
The reversible gauge sector gives a well-defined resource preorder at fixed
resource density. A richer irreversible conversion theory can be added by
including physically specified lossy, coarse-graining, or truncation maps and
requiring them to be contractive for $r$. Such maps should not be postulated as
an abstract transpose-covariant superchannel class, because the earlier
formulation had no laboratory meaning. The present paper therefore does not
claim a complete classification of all irreversible laboratory conversions.
\end{remark}

% ================================================================
\section{Numerical validation}
\label{sec:numerics}

The numerical quantity \(\PhiC\) is obtained directly from the finite matrix;
no fixed-time approximation enters. For simple spectra, there are two
numerically equivalent routes:
\begin{equation}
\kappa_\lambda=\frac{1}{|\ell_\lambda^\dagger r_\lambda|}
\end{equation}
with unit-normalized left/right eigenvectors, or
\begin{equation}
\kappa_\lambda=\norm{P_\lambda}
\end{equation}
from the corresponding spectral projector. A Schur-based implementation is
preferred near clustered eigenvalues; arbitrary-precision arithmetic may be
necessary at large $L$ because NHSE condition numbers can themselves grow
exponentially and thereby destabilize standard double-precision
calculations \cite{FengEtAl2025}.

A complete numerical validation should contain two independent checks.

\subsection{Hatano--Nelson scaling}
For fixed $g\ne0$, one should compute \(\PhiC(H_{g,L}^{\rm OBC})\) and
\(\PhiC(H_{g,L}^{\rm PBC})\) over increasing $L$. The required signatures are
\begin{equation}
\frac{\PhiC(H_{g,L}^{\rm OBC})}{L}\longrightarrow |g|,
\qquad
\PhiC(H_{g,L}^{\rm PBC})=0,
\end{equation}
with the stronger finite-size relation
\begin{equation}
\PhiC(H_{g,L}^{\rm OBC})-|g|(L-1)+\log(L+1)=O_g(1).
\end{equation}
A log-linear plot of the geometric mean projector condition number is useful,
but the primary fit should be performed on the exact \(\PhiC\) defined in
Eq.~\eqref{eq:Phi}, not on an eigenvector-matrix condition number with an
unspecified normalization convention.

\subsection{Multiband GBZ check}
A second test should use a genuinely matrix-valued finite-range symbol with a
nontrivial GBZ branch structure, for example a non-Hermitian SSH-type chain.
For each size $L$, one should independently compute the OBC spectral
projectors, form \(\PhiC\), and compare \(\PhiC/L\) with the numerical GBZ
integral in Eq.~\eqref{eq:GBZ-resource}. The purpose of this test is not to
fit a new exponent to the same HN model but to check the multiband content of
the theorem.

\begin{remark}[Numerical evidence versus theorem]
A numerical agreement can test the finite-size implementation and the
GBZ correspondence, but it cannot replace the regularity assumptions in
Assumption~\ref{ass:GBZ}. In particular, a calculation at one fixed size
cannot establish an extensive thermodynamic limit.
\end{remark}

% ================================================================
\section{Relation to existing non-normality diagnostics}
\label{sec:prior}

Condition numbers and other scalar measures of non-normality have already
been shown to grow macroscopically in NHSE phases, and the same enhancement
controls perturbation sensitivity and numerical instability
\cite{NakaiEtAl2024,FengEtAl2025}. The contribution here is not the invention
of condition numbers as NHSE diagnostics. Instead, we package the
\emph{boundary excess} of a unitary-invariant biorthogonal conditioning
functional into a resource object with explicit free transformations and a
thermodynamic rate.

This distinction also clarifies what the theory does not claim. The resource
is not a universal order parameter for every non-Hermitian phase. Generic
non-normality unrelated to boundary sensitivity may contribute to
\(\PhiC(H_L^{\rm OBC})\) and \(\PhiC(H_L^{\rm PBC})\) alike and cancel from the
boundary excess. Conversely, critical or scale-free NHSE can require a
different finite-size normalization because its localization length scales
with system size \cite{LiSunYangLi}.

% ================================================================
\section{Conclusion}
\label{sec:conclusion}

We formulated a boundary-sensitive operator resource theory for the
non-Hermitian skin effect. The essential structural choice is the object
class: the theory acts on the pair of OBC and PBC realizations of the same
bulk rule. The scalar resource is the open-boundary excess of the mean
logarithmic norm of spectral projectors. This choice gives three properties
that a single-boundary, basis-dependent similarity cost could not provide
simultaneously: exact invariance under global unitary relabeling, stability
under uniformly conditioned local recoupling, and an extensive OBC--PBC
contrast.

For Hatano--Nelson, the resource is analytically controlled:
\(\mathcal R_L=|g|L-\log L+O(1)\), with zero PBC contribution and thermodynamic
density \(r=|g|\). For a regular finite-range translation-invariant class, the
same density is given by the average GBZ decay exponent. The latter statement
is intentionally conditional: the theory treats finite-skin-depth NHSE as an
extensive resource and does not force scale-free or critical limits into the
same scaling law.

The principal remaining empirical task is therefore straightforward rather
than structural: numerical evaluation of the defined projector condition
number for the HN family and for at least one genuinely multiband GBZ model,
with sufficient numerical precision to resolve exponentially ill-conditioned
OBC eigenvectors. The theoretical resource itself does not rely on fixed-time
dynamics.

% ================================================================
\begin{acknowledgments}
\end{acknowledgments}

% ================================================================
\begin{thebibliography}{99}

\bibitem{ChitambarGour}
E.~Chitambar and G.~Gour,
\textit{Quantum resource theories},
Rev. Mod. Phys. \textbf{91}, 025001 (2019).

\bibitem{LiuYuan}
Y.~Liu and X.~Yuan,
\textit{Operational resource theory of quantum channels},
Phys. Rev. Research \textbf{2}, 012035(R) (2020).

\bibitem{ChiribellaSupermap}
G.~Chiribella, G.~M.~D'Ariano, and P.~Perinotti,
\textit{Transforming quantum operations: quantum supermaps},
Eur. Phys. Lett. \textbf{83}, 30004 (2008).

\bibitem{HatanoNelson}
N.~Hatano and D.~R.~Nelson,
\textit{Localization transitions in non-Hermitian quantum mechanics},
Phys. Rev. Lett. \textbf{77}, 570 (1996).

\bibitem{YaoWang}
S.~Yao and Z.~Wang,
\textit{Edge states and topological invariants of non-Hermitian systems},
Phys. Rev. Lett. \textbf{121}, 086803 (2018).

\bibitem{OkumaEtAl2020}
N.~Okuma, K.~Kawabata, T.~Shiozaki, and M.~Sato,
\textit{Topological origin of non-Hermitian skin effects},
Phys. Rev. Lett. \textbf{124}, 086801 (2020).

\bibitem{NakaiEtAl2024}
Y.~O.~Nakai, N.~Okuma, D.~Nakamura, K.~Shimomura, and M.~Sato,
\textit{Topological enhancement of nonnormality in non-Hermitian skin effects},
Phys. Rev. B \textbf{109}, 144203 (2024).

\bibitem{FengEtAl2025}
X.~Feng, S.~Liu, S.-X.~Zhang, and S.~Chen,
\textit{Numerical instability of non-Hermitian Hamiltonian evolution},
Phys. Rev. B \textbf{111}, 224310 (2025).

\bibitem{LiSunYangLi}
W.~Li, Z.~Sun, Z.~Yang, and F.~Li,
\textit{Universal scalefree non-Hermitian skin effect near the Bloch point},
Phys. Rev. B \textbf{109}, 035119 (2024).

\bibitem{ChienLiuNakazatoTam}
M.-T.~Chien, J.~Liu, H.~Nakazato, and T.-Y.~Tam,
\textit{Toeplitz matrices are unitarily similar to symmetric matrices},
Linear Multilinear Algebra \textbf{65}, 2131--2144 (2017).

\bibitem{TrefethenEmbree}
L.~N.~Trefethen and M.~Embree,
\textit{Spectra and Pseudospectra: The Behavior of Nonnormal Matrices and
Operators} (Princeton University Press, Princeton, NJ, 2005).

\end{thebibliography}
\end{document}
